% SPDX-License-Identifier: GPL-3.0-or-later
% Copyright (C) 2026 Alessandro Greco
\chapter{Установка и torrc}\label{cap:17-torrc}
\aggiornatoal{23 сентября 2026 г.}{tor 0.4.9.13; руководство \texttt{tor(1)} для той же версии}

\begin{obiettivi}
К концу этой главы вы сможете:
\begin{itemize}[nosep]
  \item выбирать подходящий способ установки Tor в зависимости от задачи;
  \item читать и писать файл \file{torrc} и знать, где искать документацию по каждой опции;
  \item настраивать tor как клиент, включая SOCKS-порты, изоляцию потоков и защищённый порт управления;
  \item проверять конфигурацию до её использования и диагностировать проблемы по журналам.
\end{itemize}
\end{obiettivi}

С этой главы начинается практическая часть книги. В предыдущих главах мы разбирали, как работает протокол; теперь научимся использовать и настраивать программу, которая его реализует. Все упомянутые здесь параметры описаны в руководстве tor (\cmd{man tor}), которое является авторитетным источником: значения по умолчанию могут меняться между версиями.

% ---------------------------------------------------------------------------
\section{Какой вариант Tor устанавливать?}\label{sec:torrc-installare}

\begin{description}
  \item[Tor Browser] Для просмотра веб-страниц это почти всегда правильный выбор. Он включает предварительно настроенный tor и браузер, модифицированный так, чтобы не раскрывать сведения, способные свести на нет защиту сети (\cref{cap:23-browser}). Напротив, использование tor с обычным браузером оставляет множество возможностей для идентификации пользователя.
  \item[Пакет \texttt{tor}] Подходит для использования Tor с другими приложениями, а также для запуска реле или onion-сервиса. Дистрибутивы Linux обычно включают этот пакет, хотя не всегда последней версии. Tor Project поддерживает собственный репозиторий пакетов для Debian и Ubuntu и распространяет \emph{Tor Expert Bundle}, содержащий tor и подключаемые транспорты, для Windows и macOS.
  \item[Сборка из исходного кода] Как в лабораторной среде этой книги (\cref{sec:rp-avvio}); полезна для изучения tor или получения строго определённой версии.
  \item[Arti] Клиент на Rust (\cref{cap:30-arti}), который становится всё более зрелым, но настраивается иначе.
\end{description}

Какой бы способ вы ни выбрали, загружайте программу только с официального сайта и проверяйте её подпись. Версия tor, модифицированная атакующим, свела бы его защитные свойства к нулю.

% ---------------------------------------------------------------------------
\section{Файл torrc}\label{sec:torrc-formato}

tor читает конфигурацию из текстового файла, обычно \file{/etc/tor/torrc}, если работает как системная служба. Формат прост: одна опция на строку, за которой следуют её значения. Строки, начинающиеся с \texttt{\#}, являются комментариями. Некоторые опции могут встречаться несколько раз, например несколько SOCKS-портов. Параметры также можно передавать в командной строке с двумя ведущими дефисами:
\begin{lstlisting}[language=bash,numbers=none]
$ tor --SocksPort 9150 --DataDirectory /tmp/tor-test
\end{lstlisting}
Руководство описывает сотни опций; команда \cmd{tor --list-torrc-options} выводит их все. Дистрибутивы могут задавать второй файл значений по умолчанию через \texttt{--defaults-torrc}. Файл, указанный параметром \texttt{-f}, имеет приоритет.

\subsection{Проверка перед использованием}

Три команды позволяют избежать многих неожиданностей:
\begin{lstlisting}[language=bash,numbers=none]
$ tor --verify-config -f /etc/tor/torrc     # конфигурация корректна?
Configuration was valid
$ tor --dump-config short -f /etc/tor/torrc # что отличается от значений по умолчанию?
$ tor --list-modules                        # какие модули собраны?
relay: yes
dirauth: yes
dircache: yes
pow: yes
\end{lstlisting}
Как мы видели в \cref{cap:15-dos}, последняя команда, например, показывает, доступно ли доказательство выполнения работы (PoW) для onion-сервисов.

% ---------------------------------------------------------------------------
\section{Конфигурация клиента}\label{sec:torrc-client}

Клиенту требуется очень мало настроек. Фактически минимальный конфигурационный файл пуст: по умолчанию tor открывает SOCKS-порт 9050 на localhost. Ниже приведён более явный пример с комментариями для системной службы tor, у которой есть права на каталоги \file{/var/lib/tor} и \file{/var/log/tor}; он не рассчитан на запуск от произвольного непривилегированного пользователя:
\begin{lstlisting}[language=torrc]
# Где tor хранит состояние, ключи и кеш
DataDirectory /var/lib/tor

# Порт SOCKS для приложений (только localhost!)
SocksPort 127.0.0.1:9050

# Порт управления с cookie-аутентификацией (см. главу 18)
ControlPort 127.0.0.1:9051
CookieAuthentication 1

# Журнал с уровнем notice
Log notice file /var/log/tor/notice.log
\end{lstlisting}

\begin{figure}[htbp]
\centering
% SPDX-License-Identifier: GPL-3.0-or-later
% Copyright (C) 2026 Alessandro Greco
%% Альтернативное описание для EPUB.
\altfig{Пример архитектуры tor: браузер и curl подключаются к SocksPort 9050 на localhost, а Stem и Nyx используют ControlPort 9051, только если он включён и защищён аутентификацией по cookie или паролю. Локальные протоколы SOCKS и управления сами по себе не шифруют трафик, поэтому эти порты нельзя открывать для доступа из сети. DataDirectory — каталог с ключами, состоянием и кешем, а не порт. Реле, включая мост, может открыть для других узлов ORPort 443; число 443 приведено лишь для примера.}
% Пример портов и каталога данных процесса tor
\begin{tikzpicture}[font=\sffamily\scriptsize,
  porta/.style={draw,rounded corners=1pt,fill=torviola!10,minimum height=0.55cm,minimum width=2.9cm,font=\ttfamily\scriptsize},
  dir/.style={draw=torgrigio,rounded corners=1pt,fill=torgrigio!10,minimum height=0.55cm,minimum width=2.9cm,font=\ttfamily\scriptsize},
  app/.style={draw=torgrigio,fill=torgrigio!8,rounded corners=2pt,minimum height=0.6cm,minimum width=2.1cm}]
  \node[draw=torviola,fill=torviola!5,rounded corners=4pt,minimum width=3.6cm,minimum height=4.3cm] (tor) at (4.6,0) {};
  \node[font=\sffamily\small,text=torviola] at (4.6,1.85) {\textbf{процесс tor}};
  \node[porta] (socks) at (4.6,1.0) {SocksPort 9050};
  \node[porta] (ctrl) at (4.6,0.2) {ControlPort 9051};
  \node[porta] (or) at (4.6,-0.6) {ORPort 443};
  \node[dir] (dd) at (4.6,-1.5) {DataDirectory};
  \node[app] (br) at (0.6,1.0) {браузер, curl};
  \node[app] (st) at (0.6,0.2) {Stem, Nyx};
  \node[app,fill=torverde!12,draw=torverde!60!black] (rete) at (9.0,-0.6) {сеть Tor};
  \draw[->] (br) -- (socks); \draw[<->] (st) -- (ctrl); \draw[<->] (or) -- (rete);
  \node[anchor=west,font=\sffamily\tiny,align=left,text=torgrigio] at (6.6,1.0) {только localhost:\\SOCKS без шифрования};
  \node[anchor=west,font=\sffamily\tiny,align=left,text=torgrigio] at (6.6,0.2) {если включён:\\cookie или пароль};
  \node[anchor=west,font=\sffamily\tiny,align=left,text=torgrigio] at (6.6,-1.5) {ключи, состояние guard,\\кеш консенсуса};
  \node[font=\sffamily\tiny,text=torgrigio,align=center] at (4.6,-2.5) {ORPort 443 — пример\\для реле и мостов};
\end{tikzpicture}

\caption{Пример портов и каталога процесса tor. SOCKS на порту 9050 обслуживает локальные приложения; ControlPort 9051 не включён по умолчанию и при включении требует аутентификации. Эти порты нельзя открывать для внешней сети. DataDirectory — каталог, а не порт. ORPort 443 — лишь пример настройки реле или моста, а не обязательный номер порта.}
\label{fig:porte}
\end{figure}

\begin{attenzione}[Не публикуйте SOCKS-порт]
Руководство формулирует это однозначно: протокол SOCKS не обеспечивает ни шифрования, ни аутентификации. \opt{SocksPort}, слушающий адрес, доступный из сети, превращает компьютер в открытый прокси и раскрывает запросы приложений в незашифрованном виде в локальной сети.
\end{attenzione}

\subsection{Изоляция потоков}\label{sec:torrc-isolamento}

Несколько потоков могут использовать одну и ту же цепь. Это эффективно, но позволяет выходному реле (exit) связать разные действия одного пользователя. Если личная электронная почта и псевдонимный форум используют одну цепь, выходное реле может определить, что эти действия исходят от одного и того же человека. Поэтому tor поддерживает \textbf{изоляцию потоков}: правила, определяющие, какие потоки могут совместно использовать цепь. Эти правила задаются как параметры \opt{SocksPort}:
\begin{description}
  \item[\texttt{IsolateClientAddr}] Потоки от разных клиентских адресов не могут использовать одни и те же цепи; включено по умолчанию.
  \item[\texttt{IsolateSOCKSAuth}] Потоки с разными SOCKS-учётными данными не могут совместно использовать цепи; включено по умолчанию.
  \item[\texttt{IsolateDestAddr}, \texttt{IsolateDestPort}] Изоляция по адресу или порту назначения.
  \item[\texttt{IsolateClientProtocol}] Изоляция по протоколу клиента (SOCKS4, SOCKS5, DNS\ldots).
\end{description}
Разные SOCKS-порты по умолчанию также изолированы друг от друга: каждому прослушивающему порту назначается отдельная \emph{группа сеансов (session group)}. Опция \opt{SessionGroup} позволяет изменить это поведение и разрешить портам совместно использовать цепи. Наиболее распространённый механизм — изоляция по SOCKS-аутентификации: Tor Browser отправляет разные SOCKS-учётные данные для каждого посещаемого сайта, получая отдельную цепь для каждого сайта. Мы увидим это в \cref{cap:18-stem}.

Цепь не используется бесконечно. После истечения \opt{MaxCircuitDirtiness} — по умолчанию 10 минут, отсчитываемых от первого использования — цепь перестаёт принимать новые потоки, и tor строит новую.

\subsection{Выбор реле: соблюдайте осторожность}

Опции \opt{EntryNodes}, \opt{ExitNodes} и \opt{ExcludeNodes} позволяют принудительно выбирать или исключать реле по отпечатку, псевдониму, адресу или коду страны. \opt{StrictNodes} делает требования \opt{ExcludeNodes} обязательными даже в тех случаях, когда это нарушает работоспособность. Эти инструменты полезны для экспериментов и особых сценариев. Однако использовать их только для того, чтобы \enquote{чувствовать себя безопаснее}, обычно вредно: сужение выбора реле уменьшает разнообразие маршрутов, делает пользователя более заметным и подрывает защиту, обеспечиваемую взвешенным выбором маршрута и входных реле (guard) (\cref{cap:13-percorso}).

\subsection{Другие порты}

\begin{description}
  \item[\opt{DNSPort}] tor принимает DNS-запросы и разрешает имена через сеть (\cref{sec:stream-dns}).
  \item[\opt{TransPort}] Принимает \emph{прозрачные} TCP-соединения, перенаправленные системным межсетевым экраном. Такие системы, как Whonix, используют этот механизм для маршрутизации трафика приложений через Tor (\cref{cap:23-browser}).
  \item[\opt{AutomapHostsOnResolve}] Назначает виртуальные адреса именам \texttt{.onion}, запрошенным через DNS, что позволяет приложениям без поддержки SOCKS обращаться к onion-сервисам через \opt{TransPort}.
\end{description}

% ---------------------------------------------------------------------------
\section{Порт управления}\label{sec:torrc-controllo}

\opt{ControlPort} позволяет другим программам управлять работающим tor: читать его состояние, запрашивать новые цепи или динамически создавать onion-сервисы. Это чрезвычайно мощный интерфейс, поэтому он должен быть защищён. Доступны два метода аутентификации:
\begin{description}
  \item[Файл cookie] При \opt{CookieAuthentication} 1 tor записывает в каталог данных файл, содержащий 32 случайных байта. Любой, кто может прочитать этот файл, может пройти аутентификацию. Следовательно, безопасность зависит от прав доступа к файлу.
  \item[Пароль] При \opt{HashedControlPassword} tor хранит хэш пароля, созданный командой \cmd{tor --hash-password}:
\begin{lstlisting}[language=bash,numbers=none]
$ tor --hash-password test
16:458528FE0C0AF36260B5071446823FB789212185D7EE7F615BC120D55D
\end{lstlisting}
Показанный хеш соответствует демонстрационному паролю \texttt{test} и вычислен для фиксированной соли. При настоящем вызове команда выбирает случайную соль, поэтому результат обычно будет другим даже для того же пароля. Нельзя повторно использовать пароль или хеш из примера для защиты реального порта.
\end{description}
Как и порт SOCKS, порт управления никогда не должен быть доступен из сети. Протокол управления рассматривается в \cref{cap:18-stem}.

% ---------------------------------------------------------------------------
\section{Журналы}\label{sec:torrc-log}

tor записывает события с несколькими уровнями детализации: \texttt{err}, \texttt{warn}, \texttt{notice}, \texttt{info} и \texttt{debug}. Для обычной эксплуатации достаточно уровня \texttt{notice}. Уровни \texttt{info} и \texttt{debug} полезны при исследовании протокола, как в лабораторных работах, но создают большой объём данных. Опция \opt{SafeLogging}, включённая по умолчанию, заменяет адреса и назначения в журналах на \texttt{[scrubbed]}. Забытый файл журнала не должен превращаться в историю действий пользователя. В лабораторной сети эта защита намеренно отключена (\texttt{SafeLogging 0}), чтобы можно было исследовать все детали.

Наиболее полезны при запуске сообщения \emph{bootstrap}. Они показывают этапы подключения к сети и заканчиваются строкой \texttt{Bootstrapped 100\% (done): Done}. Если процесс останавливается на определённом проценте, соответствующий этап подсказывает направление диагностики. Например, зависание при подключении к первому реле обычно указывает на цензуру или межсетевой экран и говорит в пользу использования мостов (\cref{cap:16-censura}).

% ---------------------------------------------------------------------------
\section*{Итоги}
\begin{itemize}
  \item Для просмотра веб-страниц используйте Tor Browser; пакет \texttt{tor} предназначен для других приложений, реле и onion-сервисов.
  \item Файл \file{torrc} содержит одну опцию на строку. Руководство является авторитетным источником, а \texttt{--verify-config} проверяет конфигурацию.
  \item Порты SOCKS и управления должны оставаться на localhost. Защищайте порт управления файлом cookie или паролем.
  \item Изоляция потоков не позволяет несвязанным действиям совместно использовать одну цепь. Tor Browser достигает этого с помощью разных SOCKS-учётных данных для каждого сайта.
  \item Принудительный выбор или исключение реле обычно снижает безопасность; \opt{SafeLogging} защищает содержимое журналов.
\end{itemize}

% ---------------------------------------------------------------------------
\section*{Упражнения}

\begin{esercizio}[1]\label{es:17-verifica}
В лабораторном контейнере создайте \file{torrc} с опечаткой в имени одной из опций и проверьте его с помощью \texttt{--verify-config}. Какое сообщение вы получите? Затем выполните \texttt{--dump-config short} для файла \file{torrc} клиента лабораторной сети.
\end{esercizio}

\begin{esercizio}[2]\label{es:17-isolamento}
Пользователь подключает личный почтовый клиент под своей реальной идентичностью и псевдонимный форум через один и тот же \opt{SocksPort}. Что может вывести из этого выходное реле? Предложите две конфигурации, позволяющие избежать проблемы.
\end{esercizio}

\begin{esercizio}[2]\label{es:17-exitnodes}
Почему \texttt{ExitNodes \{it\}}, то есть использование только итальянских выходных реле, снизило бы анонимность итальянского пользователя, посещающего итальянский веб-сайт? Обоснуйте ответ в терминах модели угроз из \cref{cap:04-modello-minaccia}.
\end{esercizio}

% ---------------------------------------------------------------------------
\section*{Дополнительная литература}
\begin{itemize}
  \item Руководство tor, \cmd{man tor}, и \file{doc/man/tor.1.txt} в дереве исходного кода.
  \item Страницы поддержки Tor Project по установке и репозиториям пакетов: \url{https://support.torproject.org/}.
\end{itemize}
